\begin{textsample}{POS Dim 7 – human – Score 4.00 – rj/rj\_ef\_linguagens\_ed\_fis\_3\_p11.txt}  \label{ex:f7_pos_009}
Como dito anteriormente, o DOC-RJ propõe as \textbf{seguintes} Unidades Temáticas conforme previsto na BNCC ( 2017, p.212 ) : Brincadeiras e Jogos ; Esportes ; Ginásticas ; Danças ; Lutas e Práticas Corporais de Aventura.
Danças e Lutas são organizadas conforme a ocorrência social dessas práticas ( \textbf{local}, \textbf{regional}, nacional e mundial ), já as Ginásticas apoiam -se em sua diversidade e características ( Geral, de Condicionamento Físico e de Conscientização Corporal ), nos Esportes a \textbf{abordagem} considera o modelo de classificação ( Marca, Precisão, Técnico-combinatório, Rede/Quadra dividida ou Parede de rebote, Campo e taco, Invasão ou Territorial e Combate ) e nas Práticas Corporais de Aventura destaca -se uma exposição urbana e \textbf{natural}.

% matched lemmas: abordagem, local, natural, regional, seguinte
\end{textsample}
